% !TeX root = ../../MainBook.tex
\section{线性差分方程组与状态变量分析}\label{sec:Difference_Equation}

\subsection{线性差分方程组与常系数线性差分方程组}\label{sub:linear_difference_equation}

\begin{definition}[一阶线性差分方程组]
    设$A[n]$为$N\times N$的实矩阵值函数，$\mathbf{f}[n]$为$N$维实向量值函数，则形如
    \begin{equation}
        \mathbf{y}[n+1]=A[n]\mathbf{y}[n]+\mathbf{f}[n]
    \end{equation}
    的方程组称为\textbf{一阶线性（前向）差分方程组}，其中$\mathbf{f}[n]$为\textbf{非齐次项}。

    特别地，如果$\mathbf{f}[n]=\mathbf{0}$，则称之为\textbf{齐次线性差分方程组}；如果$A[n]$为常数矩阵，则称为\textbf{常系数线性差分方程组}。

    将方程的个数$N$称为方程组的\textbf{阶数}。
\end{definition}

不难看出，类似于\refsub{compare_to_high_order}中微分方程组与微分方程的关系，\textbf{差分方程组也能与同阶的（多个）差分方程相互转化}。考虑前向差分方程组，是因为方程组不涉及因果性的问题，而前向差分方程组是更简单、常见的。

任给正整数$k$，可以通过迭代法求出
\begin{align}
    \mathbf{y}[k]&=A[k-1]\mathbf{y}[k-1]+\mathbf{f}[k-1] \nonumber \\
    &=A[k-1]\paren{A[k-2]\mathbf{y}[k-2]+\mathbf{f}[k-2]}+\mathbf{f}[k-1] \nonumber \\
    &=A[k-1]A[k-2]\mathbf{y}[k-2]+A[k-1]\mathbf{f}[k-2]+\mathbf{f}[k-1] \nonumber \\
    &\ \ \vdots \nonumber \\
    &=A[k-1]A[k-2]\cdots A[0]\mathbf{y}[0]+\sum_{i=0}^{k-1}A[k-1]A[k-2]\cdots A[i+1]\mathbf{f}[i].
\end{align}
因此，给定初始值$\mathbf{y}[0]$，就可以\textbf{唯一确定}任意时刻的$\mathbf{y}[k]$。特别地，如果$A[n]$为常数矩阵，则
\begin{equation}
    \mathbf{y}[k]=A^k\mathbf{y}[0]+\sum_{i=0}^{k-1}A^{k-1-i}\mathbf{f}[i].
\end{equation}

利用解的存在性与唯一性，可以仿照\cref{thm:n_order_ODE}的方法进一步说明：
\begin{theorem}
    $N$阶线性差分方程组的解空间$S$构成$N$维实线性空间。
\end{theorem}
\begin{proof}
    线性性是显然的。我们构造一个同构映射来说明解空间的维数是$N$。

    定义$\mathbf{H}:\mathbb{R}^N\to S,\ \mathbf{y}_0\mapsto\mathbf{y}[n]$，其中$\mathbf{y}[n]$是满足初始条件$\mathbf{y}[0]=\mathbf{y}_0$的解。下验证$\mathbf{H}$是一个线性空间同构映射：

    首先，$\mathbf{H}$是线性的。$\forall \mathbf{y}_0,\mathbf{z}_0\in\mathbb{R}^N,\ \alpha,\beta\in\mathbb{R}$，$\alpha\mathbf{H}(\mathbf{y}_0)+\beta\mathbf{H}(\mathbf{z}_0)$是满足初始条件$\alpha\mathbf{y}_0+\beta\mathbf{z}_0$的解，根据解的唯一性，
    \begin{equation}
        \mathbf{H}(\alpha\mathbf{y}_0+\beta\mathbf{z}_0)=\alpha\mathbf{H}(\mathbf{y}_0)+\beta\mathbf{H}(\mathbf{z}_0).
    \end{equation}

    其次，$\mathbf{H}$是单射。设$\mathbf{y}_0\neq\mathbf{z}_0$，它们对应的解分别为$\mathbf{y}=\mathbf{H}(\mathbf{y}_0)$和$\mathbf{z}=\mathbf{H}(\mathbf{z}_0)$，则$\mathbf{y}[0]\neq\mathbf{z}[0]$，因此$\mathbf{y}\neq\mathbf{z}$。

    最后，$\mathbf{H}$是满射。对于任意$\mathbf{y}[n]\in S$，$\mathbf{y}_0=\mathbf{y}[0]$满足$\mathbf{H}(\mathbf{y}_0)=\mathbf{y}[n]$。

    综上，$\mathbf{H}$是一个线性空间的同构映射，解空间的维数为$\dim\mathbb{R}^N=N$。
\end{proof}

由于高阶线性差分方程总能转化为同阶的线性差分方程组，因此高阶线性差分方程的解空间也是一个维数为$N$的线性空间，这就解释了为什么上一节中我们总是给出$N$个初始值来唯一确定一个高阶线性差分方程的解。

前面说明解的存在性与唯一性时，已经给出了通解的表达式，只是在$A[n]$不为常数值矩阵时，没有什么比较好的办法能够简化解的计算，除非$A[n]$具有某些特殊的结构，例如对角矩阵、上三角矩阵等，但这个问题在常系数线性差分方程组中是易于解决的。

对于常系数线性差分方程组，我们也已经给出了通解的表达式：
\begin{equation}
    \mathbf{y}[k]=A^k\mathbf{y}[0]+\sum_{i=0}^{k-1}A^{k-1-i}\mathbf{f}[i].
\end{equation}
因此，求解常系数线性差分方程组的关键在于求出$A^k$的表达式。在多数问题中，我们只需要求出有限的几个$\mathbf{y}[k]$，除了直接迭代，通常有两种快速算法：Jordan分解算法和Cayley-Hamilton定理递推算法。

假设已知矩阵$A$的Jordan分解$A=PJP^{-1}$，其中$J$是一个块对角矩阵，就可以很容易地求出$J^k$的表达式，从而
\begin{equation}
    A^k=(PJP^{-1})^k=PJ^k P^{-1}.
\end{equation}

Cayley-Hamilton定理（见\refsub{eigenvector}）的内容是，设矩阵$A$的特征多项式为
\begin{equation}
    p(\lambda)=\det(\lambda I-A)=\sum_{i=0}^{N}c_i\lambda^i,
\end{equation}
则
\begin{equation}
    p(A)=\sum_{i=0}^{N}c_i A^i=0.
\end{equation}
因此，$A^N$可以表示为$A^0,A^1,\cdots,A^{N-1}$的线性组合，即
\begin{equation}
    A^N=-\sum_{i=0}^{N-1}c_i A^i.
\end{equation}
这样，对于任意的$k\geq N$，我们都可以将$A^k$表示为$A^0=I,A^1,\cdots,A^{N-1}$的线性组合，并递归地求出线性组合的系数，从而求出$\mathbf{y}[k]$的表达式。

下面给出并比较三种算法的\textbf{时间复杂度}。

（1）直接迭代

直接利用递推式进行计算：
\begin{equation}
    \mathbf{y}[k+1]=A\mathbf{y}[k]+\mathbf{f}[k].
\end{equation}
每一步都需要计算复杂度为$O(N^2)$的矩阵乘法，因此递归地求取$\mathbf{y}[k]$的复杂度为$O(k N^2)$。

这一方法没有预处理的时间成本，并且同样适用于$A[n]$为非恒定矩阵的情况，但当$k$较大时，计算效率较低。

（2）Jordan分解算法

在这种算法中，我们首先需要将$A$做Jordan分解，得到$A=PJP^{-1}$，这个过程包括求特征值和特征向量，以及求矩阵$P$的逆，复杂度为$O(N^3)$。然后，我们需要计算$J^k$，由于$J$是一个块对角矩阵，$J^k$的每个块都有已知的解析表达式，因此$J^k$的计算复杂度为$O(N^2)$。因此，Jordan分解算法的预处理复杂度为$O(N^3)$。

对于齐次方程，每求一个$\mathbf{y}[k]$，都要计算矩阵乘法$A\mathbf{y}[k]=PJ^k P^{-1}\mathbf{y}[0]$，复杂度为$O(N^3)$，因此每求一个$\mathbf{y}[k]$的复杂度为$O(N^3)$。

对于非齐次方程，由于需要计算$\sum_{i=0}^{k-1}A^{k-1-i}\mathbf{f}[i]$，每次矩阵乘法$A^{k-1-i}\mathbf{f}[i]$的复杂度为$O(N^2)$，因此每求一个$\mathbf{y}[k]$的复杂度为$O(N^3)+O(k N^2)$。

（3）Cayley-Hamilton定理递推算法

为了计算矩阵的特征多项式，通常使用Hessenberg分解算法，复杂度为$O(N^3)$；通过矩阵乘法计算$A^i,\ i=2,3,\cdots,N-1$，复杂度为$O(N^3)$。在预处理时，还需计算$\mathbf{y}[0],\mathbf{y}[1],\cdots,\mathbf{y}[N-1]$，以免后面将$A^k$表示为$A^0=I,A^1,\cdots,A^{N-1}$的线性组合之后进行重复的矩阵乘法运算，这个过程的复杂度为$O(N^3)$。因此这种算法的预处理复杂度为$O(N^3)$。

对于齐次方程，$\mathbf{y}[k]$可以表示为$\mathbf{y}[0],\mathbf{y}[1],\cdots,\mathbf{y}[N-1]$的线性组合，系数的求解可以通过多项式快速幂算法来实现，复杂度为$O(N^2\log k)$；矩阵乘法的复杂度为$O(N^2)$。因此每求一个$\mathbf{y}[k]$的复杂度为$O(N^2\log k)$。

对于非齐次方程，由于需要计算$\sum_{i=0}^{k-1}A^{k-1-i}\mathbf{f}[i]$，如果每个$A^{k-1-i}$都通过多项式快速幂算法来计算，那么每求一个$\mathbf{y}[k]$的复杂度为$O(k N^2\log k)$，此时复杂度比直接迭代还要大。

将上述三种算法在齐次/非齐次方程下的复杂度汇总如下：
\begin{table}[H]
    \centering
    \begin{tabular}{|c|c|c|}
        \hline
        算法                         & 齐次方程                       & 非齐次方程                         \\
        \hline
        直接迭代                     & $O(k N^2)$                     & $O(k N^2)$                         \\
        \hline
        Jordan分解算法               & $O(N^3)$                       & $O(N^3)+O(k N^2)$                  \\
        \hline
        Cayley-Hamilton定理递推算法  & $O(N^3)+O(N^2 \log k)$         & $O(N^3)+O(k N^2 \log k)$           \\
        \hline
    \end{tabular}
\end{table}

可见，当只需求解较小的$k$或方程含有非齐次项时，直接迭代算法的效率更高；当$k\gg N$且方程齐次时，Jordan分解算法效率最高。但是，Jordan分解算法的复杂度中$N^3$的系数较大，有时使用Cayley-Hamilton定理递推算法会更快一些。

\subsection{解矩阵}\label{sub:fundamental_solution_matrix}

类比微分方程组，定义：
\begin{definition}[解矩阵]
    常系数线性差分方程组
    \begin{equation}
        \mathbf{y}[n+1]=A\mathbf{y}[n]+\mathbf{f}[n]
    \end{equation}
    的解矩阵为
    \begin{equation}
        \phi[n]=A^n.
    \end{equation}
\end{definition}

得到这个定义的思路有两种。一种思路是考虑将齐次差分方程中的列向量$\mathbf{y}[n]$替换为一个矩阵值函数，则$\phi[n]=A^n$恰好满足方程
\begin{equation}
    \phi[n+1]=A\phi[n],\quad \phi[0]=I.
\end{equation}

另一种思路是类比微分方程。常系数线性微分方程最典型的解的形式为$\e^{\lambda t}$，而方程组的解矩阵为$\phi(t)=\e^{At}$；对于差分方程，将$\lambda^n$中的$\lambda$替换为矩阵$A$，就得到解矩阵$\phi[n]=A^n$。

解矩阵的一个重要特征是，\textbf{它的列向量是齐次差分方程的解}，因此方程的齐次解总能表示为$\phi[n]\mathbf{c}$的形式，其中$\mathbf{c}\in\mathbb{R}^N$是一个常向量。与微分方程不同的是，当$A$不可逆时，$\phi[n]$也是不可逆的，因此$\phi^{-1}[n]$并不总是存在的。

当$A$可逆时，有$\phi^{-1}[n]=A^{-n}=\phi[-n]$，它的列空间有$N$维，恰为差分方程组的解空间，此时称之为\textbf{基本解矩阵}。

事实上，对于变系数的线性差分方程组
\begin{equation}
    \mathbf{y}[n+1]=A[n]\mathbf{y}[n]+\mathbf{f}[n],
\end{equation}
我们同样可以定义解矩阵：
\begin{equation}
    \phi[n]=A[n-1]A[n-2]\cdots A[0].
\end{equation}
它满足$\phi[n+1]=A[n]\phi[n],\ \phi[0]=I$，因此其列向量确为方程组的解。但由于$A[n]$不为常数矩阵，我们通常是求不出$\phi[n]$的解析表达式的，这一点与\refsec{Equation_System}中的线性微分方程组有些类似。

解矩阵在差分方程组理论中，仍然扮演着“基本解”的角色。以非齐次项$\mathbf{f}[n]$为激励，方程组
\begin{equation}
    \mathbf{y}[n+1]=A\mathbf{y}[n]+\mathbf{f}[n],\quad \mathbf{y}_0=\mathbf{y}[0]
\end{equation}
的解$\mathbf{y}[n]$为响应，则方程描述了一个离散的因果线性时不变系统，类比\refsub{constant_coefficient_linear_equation_system}中的常系数线性微分方程组，可以猜想零状态响应可以表示为卷积和，然而，我们实际上应该使用$\phi[n-1]u[n-1]$做卷积才能得到正确的结论，我们将在下一小节中看到这样做的原因：
\begin{equation}
    \mathbf{y}_{\zs}[n]=(\phi[n-1]u[n-1])*\mathbf{f}[n]=\sum_{i=0}^{n-1}\phi[n-1-i]\mathbf{f}[i]=\sum_{i=0}^{n-1}A^{n-1-i}\mathbf{f}[i].
\end{equation}
求和范围的改变是因为我们认为激励$\mathbf{f}[n]$是在$n=0$时刻施加于系统的。这与上一小节中用迭代法得到的通解公式一致。

零输入响应是方程的齐次解。由于解矩阵的列向量是差分方程的解，零输入响应可以表示为解矩阵的列向量的线性组合：
\begin{equation}
    \mathbf{y}_{\zi}[n]=\phi[n]\mathbf{c},\quad \mathbf{c}\in\mathbb{R}^N.\label{eq:diff_sys_zi}
\end{equation}

代入初始值$\mathbf{y}[0]=\mathbf{y}_0$，可以求出$\mathbf{c}$的表达式：
\begin{equation}
    \mathbf{y}_0=\phi[0]\mathbf{c}\implies\mathbf{c}=\phi^{-1}[0]\mathbf{y}_0=\mathbf{y}_0.
\end{equation}
尽管$\phi[n]$不一定可逆，但在$n=0$时，$\phi[0]=I$。代入\refeqn{diff_sys_zi}，就得到零输入响应的表达式：
\begin{equation}
    \mathbf{y}_{\zi}[n]=\phi[n]\mathbf{c}=\phi[n]\mathbf{y}_0.
\end{equation}
最终，我们得到了常系数线性差分方程组的通解公式，它与\cref{cor:differential_equation_system}非常相似，并且与上一小节中通过迭代法得到的通解公式一致：
\begin{theorem}
    设$\phi[n]$为带有初始值的方程组
    \begin{equation}
        \mathbf{y}[n+1]=A\mathbf{y}[n]+\mathbf{f}[n],\quad \mathbf{y}_0=\mathbf{y}[0]
    \end{equation}
    的解矩阵，则方程组的通解为
    \begin{align}
        \mathbf{y}[n]&=\phi[n]\mathbf{y}_0+\sum_{i=0}^{n-1}\phi[n-1-i]\mathbf{f}[i] \nonumber \\
        &=A^n\mathbf{y}_0+\sum_{i=0}^{n-1}A^{n-1-i}\mathbf{f}[i].
    \end{align}
\end{theorem}

\subsection{常系数线性差分方程组的单边z变换解法}\label{sub:difference_equation_system_z}

单边z变换同样能够处理常系数线性差分方程组，并且能够一并处理初始值。我们首先给出离散矩阵值函数的单边z变换的定义：

\begin{definition}[离散矩阵值函数的单边z变换]
    设$A[n]$为$N\times N$的矩阵值函数，则\textbf{单边z变换}就是它的各个分量的单边z变换：
    \begin{equation}
        A(z)=\sum_{n=0}^{\infty}A[n]z^{-n},
    \end{equation}
    收敛域为各个分量的单边z变换的收敛域的交集。
\end{definition}

特别地，列向量可以视作只有一列的矩阵，因此向量值函数的单边z变换也是良定义的。

考虑如下常系数线性差分方程组
\begin{equation}
    \mathbf{y}[n+1]=A\mathbf{y}[n]+\mathbf{f}[n],\quad \mathbf{y}_0=\mathbf{y}[0],
\end{equation}
两边同时乘以$u[n]$并进行单边z变换，得到
\begin{equation}
    z\mathbf{Y}(z)-z\mathbf{y}_0=A\mathbf{Y}(z)+\mathbf{F}(z),
\end{equation}
其中$\mathbf{Y}(z)$和$\mathbf{F}(z)$分别是$\mathbf{y}[n]$和$\mathbf{f}[n]$的单边z变换。因此，
\begin{equation}
    (zI-A)\mathbf{Y}(z)=\mathbf{F}(z)+z\mathbf{y}_0.
\end{equation}
两边同时求逆z变换，即可得到方程的解：
\begin{equation}
    \mathbf{y}[n]=\mathcal{Z}^{-1}\braces{(zI-A)^{-1}(\mathbf{F}(z)+z\mathbf{y}_0)}.\label{eq:diff_sys_z_sol}
\end{equation}

令$\Phi(z)=(zI-A)^{-1}$，称之为差分方程组的\textbf{特征矩阵}或\textbf{预解矩阵}。注意到
\begin{equation}
    \Phi(z)=(zI-A)^{-1}=z^{-1}(I-z^{-1}A)^{-1}=z^{-1}\sum_{i=0}^{\infty}A^i z^{-i}=z^{-1}\mathcal{Z}\{\phi[n]u[n]\}.
\end{equation}
因此$z\Phi(z)$就是解矩阵$\phi[n]$的单边z变换（注意，这与微分方程组不同）。代入\refeqn{diff_sys_z_sol}，得到
\begin{align}
    \mathbf{y}[n]&=\mathcal{Z}^{-1}\braces{\Phi(z)(\mathbf{F}(z)+z\mathbf{y}_0)} \nonumber \\
    &=\mathcal{Z}^{-1}\braces{\Phi(z)\mathbf{F}(z)}+\mathcal{Z}^{-1}\braces{z\Phi(z)\mathbf{y}_0} \nonumber \\
    &=\mathcal{Z}^{-1}\Bigl\{z^{-1}\mathcal{Z}\{\phi[n]u[n]\}\mathbf{F}(z)\Bigr\}+\phi[n]u[n]\mathbf{y}_0.
\end{align}
利用单边z变换的移位性质以及卷积定理，得到
\begin{align}
    \mathcal{Z}^{-1}\Bigl\{z^{-1}\mathcal{Z}\{\phi[n]u[n]\}\mathbf{F}(z)\Bigr\}&=(\phi[n-1]u[n-1])*(\mathbf{f}[n]u[n]) \nonumber \\
    &=\sum_{i=0}^{n-1}\phi[n-1-i]\mathbf{f}[i] \nonumber \\
    &=\sum_{i=0}^{n-1}A^{n-1-i}\mathbf{f}[i].
\end{align}

综上，我们得到了常系数线性差分方程组的通解公式，它与之前推导的公式相一致：
\begin{equation}
    \mathbf{y}[n]=\phi[n]\mathbf{y}_0+\sum_{i=0}^{n-1}\phi[n-1-i]\mathbf{f}[i]=A^n\mathbf{y}_0+\sum_{i=0}^{n-1}A^{n-1-i}\mathbf{f}[i].
\end{equation}

\begin{example}
    求解常系数线性差分方程组在$n\in\mathbb{N}$上的解：
    \begin{equation}
        \begin{bmatrix}
            y_1[n+1] \\
            y_2[n+1]
        \end{bmatrix}=\begin{bmatrix}
            1 & 1 \\
            0 & 1
        \end{bmatrix}
        \begin{bmatrix}
            y_1[n] \\
            y_2[n]
        \end{bmatrix}+\begin{bmatrix}
            0 \\
            u[n]
        \end{bmatrix},\quad
        \begin{bmatrix}
            y_1[0] \\
            y_2[0]
        \end{bmatrix}=\begin{bmatrix}
            1 \\
            0
        \end{bmatrix}.
    \end{equation}

    解：记
    \begin{equation}
        A=\begin{bmatrix}
            1 & 1 \\
            0 & 1
        \end{bmatrix},
    \end{equation}
    两边同时乘以$u[n]$并进行单边z变换，得到
    \begin{equation}
        z\begin{bmatrix}
            Y_1(z) \\
            Y_2(z)
        \end{bmatrix}-z\begin{bmatrix}
            1 \\
            0
        \end{bmatrix}=A\begin{bmatrix}
            Y_1(z) \\
            Y_2(z)
        \end{bmatrix}+\begin{bmatrix}
            0 \\
            \frac{z}{z-1}
        \end{bmatrix},
    \end{equation}
    从而
    \begin{equation}
        (zI-A)\begin{bmatrix}
            Y_1(z) \\
            Y_2(z)
        \end{bmatrix}=\begin{bmatrix}
            z \\
            \frac{z}{z-1}
        \end{bmatrix}.
    \end{equation}
    计算特征矩阵：
    \begin{equation}
        \Phi(z)=(zI-A)^{-1}=\begin{bmatrix}
            z-1 & -1  \\
            0   & z-1
        \end{bmatrix}^{-1}=\begin{bmatrix}
            \frac{1}{z-1} & \frac{1}{(z-1)^2} \\
            0             & \frac{1}{z-1}
        \end{bmatrix}.
    \end{equation}
    因此
    \begin{align}
        \begin{bmatrix}
            Y_1(z) \\
            Y_2(z)
        \end{bmatrix}&=\Phi(z)\begin{bmatrix}
            z \\
            \frac{z}{z-1}
        \end{bmatrix} \nonumber \\
        &=\begin{bmatrix}
            \frac{1}{z-1} & \frac{1}{(z-1)^2} \\
            0             & \frac{1}{z-1}
        \end{bmatrix}\begin{bmatrix}
            z \\
            \frac{z}{z-1}
        \end{bmatrix} \nonumber \\
        &=\begin{bmatrix}
            \frac{z}{z-1}+\frac{z}{(z-1)^3} \\
            \frac{z}{(z-1)^2}
        \end{bmatrix}.
    \end{align}
    最后，选取半径大于$1$的围道$C$，求逆z变换：
    \begin{align}
        y_1[n]&=\frac{1}{2\pi\im}\oint_C \paren{\frac{z}{z-1}+\frac{z}{(z-1)^3}}z^{n-1}\,\dd z \nonumber \\
        &=\Res\paren{\frac{z^n}{z-1},1}+\Res\paren{\frac{z^n}{(z-1)^3},1} \nonumber \\
        &=\lim_{z\to 1}z^n+\lim_{z\to 1}\frac{1}{2}\frac{\dd^2}{\dd z^2}z^n \nonumber \\
        &=1+\frac{n(n+1)}{2}, \\
        y_2[n]&=\frac{1}{2\pi\im}\oint_C \frac{z}{(z-1)^2}z^{n-1}\,\dd z \nonumber \\
        &=\Res\paren{\frac{z^n}{(z-1)^2},1} \nonumber \\
        &=\lim_{z\to 1}\frac{\dd}{\dd z}z^n \nonumber \\
        &=n.
    \end{align}
\end{example}

\subsection{*离散系统的状态变量分析}\label{sub:discrete_state_variable}

对离散系统重复\refsub{state_variable_analysis}中的定义：

\begin{definition}[状态变量]
    \textbf{状态变量}(state variable)是指能够唯一确定系统状态的最小的一组变量。只要知道$n=n_0$时刻的状态变量以及$n\geq n_0$时的激励，就能够完全确定系统在$n\geq n_0$时的行为，包括状态$\mathbf{q}[n]$和响应$\mathbf{y}[n]$。

    将状态变量的个数$N$称为系统的\textbf{阶数}。
\end{definition}

\begin{definition}[状态矢量]
    \textbf{状态矢量}(state vector)是由系统的状态变量组成的向量，通常记为$\mathbf{q}[n]$。
\end{definition}

\begin{definition}[状态空间和状态轨迹]
    状态矢量所处的空间称为\textbf{状态空间}(state space)。有$n$个状态变量的系统，状态空间为$n$维实线性空间。

    随时间变化，状态向量$\mathbf{q}[n]$在状态空间中形成的轨迹称为系统的\textbf{状态轨迹}(state trajectory)。
\end{definition}

一般地，一个因果的线性时不变系统的激励$\mathbf{e}[n]$、响应$\mathbf{r}[n]$与状态变量$\mathbf{q}[n]$满足以下常系数线性差分方程组：
\begin{equation}
    \begin{aligned}
        \mathbf{q}[n+1]&=A\mathbf{q}[n]+B\mathbf{e}[n], \\
        \mathbf{r}[n]&=C\mathbf{q}[n]+D\mathbf{e}[n],
    \end{aligned}
\end{equation}
其中矩阵$A$、$B$、$C$、$D$为常数矩阵，分别称它们为系统的\textbf{状态矩阵}(state matrix)、\textbf{输入矩阵}(input matrix)、\textbf{输出矩阵}(output matrix)和\textbf{直接传递矩阵}(direct transmission matrix)。

根据前面几个小节的内容可知，已知状态变量的初始值$\mathbf{q}[0]$，关于状态变量的差分方程组的解为
\begin{equation}
    \mathbf{q}[n]=A^n\mathbf{q}[0]+\sum_{i=0}^{n-1}A^{n-1-i}B\mathbf{e}[i],
\end{equation}
因此又将方程的解矩阵$\phi[n]=A^n$称为\textbf{状态转移矩阵}(state transition matrix)，它与激励共同决定了状态矢量在状态空间中的移动情况。

考虑将系统矩阵做Jordan分解：$A=PJP^{-1}$，则状态矢量应该变为$\mathbf{w}[n]=P^{-1}\mathbf{q}[n]$，描述系统的方程变为：
\begin{equation}
    \begin{aligned}
        \mathbf{w}[n+1]&=J\mathbf{w}[n]+P^{-1}B\mathbf{e}[n], \\
        \mathbf{r}[n]&=CP\mathbf{w}[n]+D\mathbf{e}[n].
    \end{aligned}
\end{equation}

将状态变量进行上述的可逆线性变换后，\textbf{可以根据Jordan块的结构将系统拆分为若干互不耦合}（即互相独立）\textbf{的子系统的并联}，将这种操作称为状态变量的\textbf{对角化}(diagonalization)，这种分析方法称为系统的\textbf{模态分析}(modal analysis)，每个子系统称为一个\textbf{模态}(mode)。变换矩阵$P$又称为系统的\textbf{模态矩阵}(modal matrix)。

习惯上，将状态变量对角化后的系统的方程组记为
\begin{equation}
    \begin{aligned}
        \mathbf{w}[n+1]&=\hat{A}\mathbf{w}[n]+\hat{B}\mathbf{e}[n], \\
        \mathbf{r}[n]&=\hat{C}\mathbf{w}[n]+\hat{D}\mathbf{e}[n],
    \end{aligned}
\end{equation}
其中
\begin{equation}
    \hat{A}=P^{-1}AP=J,\quad \hat{B}=P^{-1}B,\quad \hat{C}=CP,\quad \hat{D}=D.
\end{equation}

特别地，如果$A$是可对角化的，则每个子系统都是一阶的，状态变量的方程为
\begin{equation}
    w_k[n+1]=\lambda_k w_k[n]+\mathbf{p}_k^T B\mathbf{e}[n],\quad k=1,2,\cdots,N,
\end{equation}
其中$\lambda_k$是$A$的第$k$个特征值，$\mathbf{p}_k^T$是$P^{-1}$的第$k$行，$\mathbf{p}_k^T B$是$\hat{B}$的第$k$行。每个子系统的状态变量$w_k[n]$的解为
\begin{equation}
    w_k[n]=\lambda_k^n w_k[0]+\sum_{i=0}^{n-1}\lambda_k^{n-1-i}\mathbf{p}_k^T B\mathbf{e}[i].
\end{equation}

在$\hat{B}$的第$k$行不全为$0$的情况下，通过适当地选取激励信号$\mathbf{e}[n]$，不论状态变量$w_k[n]$的初始值如何，都可以使得$w_k[n]$在有限的时间内达到任意的值，这表明系统的第$k$个模态是可以通过激励控制的，即该模态与输入\textbf{耦合}。系统的\textbf{可控性}(controllability)要求状态向量可以被适当的激励引向状态空间中的特定位置：

\begin{definition}[可控性]
    如果对于任意初始值$\mathbf{q}[0]$和任意目标状态$\mathbf{q}_f$，存在激励信号$\mathbf{e}[n]$能够在有限时间内使得状态矢量到达$\mathbf{q}_f$，则称系统是\textbf{可控}的。

    如果只有部分模态是可控的，则称系统是\textbf{部分可控}的。
\end{definition}

\textbf{可控性要求状态变量与输入耦合}，仅仅取决于$\hat{A}$和$\hat{B}$。如果系统矩阵$A$可对角化，则所有状态变量相互独立，\textbf{$\hat{B}$的第$k$行不全为$0$，就等价于第$k$个模态是可控的}；如果系统矩阵不可对角化，则每个模态由多个状态变量描述，子系统的系统矩阵为Jordan块：
\begin{equation}
    \begin{bmatrix}
        w_k[n+1]   \\
        w_{k+1}[n+1] \\
        \vdots     \\
        w_{k+m-1}[n+1]
    \end{bmatrix}=\begin{bmatrix}
        \lambda & 1      & 0      & \cdots & 0       \\
        0       & \lambda & 1     & \cdots & 0       \\
        \vdots  & \vdots & \vdots & \ddots & \vdots  \\
        0       & 0      & 0      & \cdots & 1       \\
        0       & 0      & 0      & \cdots & \lambda
    \end{bmatrix}\begin{bmatrix}
        w_k[n]   \\
        w_{k+1}[n] \\
        \vdots   \\
        w_{k+m-1}[n]
    \end{bmatrix}+\begin{bmatrix}
        \mathbf{p}_k^T B\mathbf{e}[n]     \\
        \mathbf{p}_{k+1}^T B\mathbf{e}[n] \\
        \vdots                            \\
        \mathbf{p}_{k+m-1}^T B\mathbf{e}[n]
    \end{bmatrix}.
\end{equation}
此时模态内的最后一个状态变量由一阶方程描述，而其余状态变量受自身、激励以及下一个状态变量的控制，因此\textbf{每个模态的可控性等价于最后一个状态变量（链尾）与输入耦合}。在\refsub{state_variable_analysis}中，我们已经展示过这一结论对二阶Jordan块对应的模态的正确性，读者可以自行构造一个类似的例子，这里不再重复。

与之类似的还有系统的\textbf{可观性}(observability)，它要求\textbf{状态变量与输出耦合}：

\begin{definition}[可观性]
    如果给定激励信号$\mathbf{e}[n]$，总能根据系统在有限时间$n_0$内的响应$\mathbf{r}[n],\ n\in[0,n_0]$唯一确定所有状态变量的起始状态$\mathbf{q}[0]$，则称系统是\textbf{可观}的。

    如果只能确定部分状态变量的起始状态，则称系统是\textbf{部分可观}的。
\end{definition}

类似地，可观性仅仅取决于$\hat{A}$和$\hat{C}$。考虑输出方程
\begin{equation}
    \begin{aligned}
        \mathbf{r}[n]&=\hat{C}\mathbf{w}[n]+\hat{D}\mathbf{e}[n], \\
        \mathbf{r}_k[n]&=\sum_{j=1}^{N}\hat{C}_{kj}w_j[n]+\sum_{j=1}^{m}\hat{D}_{kj}e_j[n].
    \end{aligned}
\end{equation}
如果系统矩阵$A$可对角化，则所有状态变量相互独立，\textbf{$\hat{C}$的第$k$行不全为$0$，就等价于第$k$个模态是可观的}，即该模态对应的子系统是可观的；如果系统矩阵不可对角化，则每个模态由多个状态变量描述，子系统的系统矩阵为Jordan块，此时模态内的后一状态变量可以由前一状态变量观测，因此\textbf{每个模态的可观性等价于第一个状态变量（链首）与输出耦合}。

下面给出离散系统的可控性与可观性的满秩判别法。

\begin{theorem}[可控性的满秩判别法]
    一个$N$阶系统是可控的，当且仅当它的\textbf{可控矩阵}(controllable matrix)
    \begin{equation}
        \mathcal{C}=\begin{bmatrix}
            B & AB & A^2B & \cdots & A^{N-1}B
        \end{bmatrix}
    \end{equation}
    是行满秩的，即$\rank\mathcal{C}=N$。
\end{theorem}
\begin{proof}
    不妨设$\mathbf{q}[0]=0$。根据状态变量的解的公式，$\forall n\in\mathbb{N}^*$，
    \begin{equation}
        \mathbf{q}[n]=\sum_{i=0}^{n-1}A^{n-1-i}B\mathbf{e}[i]=\begin{bmatrix}
            B & AB & A^2B & \cdots & A^{n-1}B
        \end{bmatrix}\begin{bmatrix}
            \mathbf e[n-1]\\
            \mathbf e[n-2] \\
            \vdots \\
            \mathbf e[0]
        \end{bmatrix}.
    \end{equation}
    矩阵右乘一个列向量，相当于对矩阵进行列变换，因此$\mathbf{q}[n]$的取值范围就等于矩阵
    \begin{equation}
        \mathcal{C}=\begin{bmatrix}
            B & AB & A^2B & \cdots & A^{n-1}B
        \end{bmatrix}
    \end{equation}
    的列空间$\mathrm{C}(\mathcal{C})$，因此$\mathbf{q}[n]$可以取$\mathbb{R}^N$中的任意值就等价于
    \begin{equation}
        \dim\mathrm{C}(\mathcal{C})=\rank(\mathcal{C})=\dim\mathbb{R}^N=N.
    \end{equation}
\end{proof}

与连续系统一样，通过对偶系统，可以由可控性的满秩判别法得到可观性的满秩判别法。
\begin{definition}[对偶系统]
    设系统的状态矩阵、输入矩阵和输出矩阵分别为$A,B,C$，则其\textbf{对偶系统}的状态矩阵、输入矩阵和输出矩阵分别为$A^T,C^T,B^T$。
\end{definition}

\begin{theorem}
    原系统的可控性等价于对偶系统的可观性，原系统的可观性等价于对偶系统的可控性。
\end{theorem}

\begin{theorem}[可观性的满秩判别法]
    一个$N$阶系统是可观的，当且仅当它的\textbf{可观矩阵}(observable matrix)
    \begin{equation}
        \mathcal{O}=\begin{bmatrix}
            C      \\
            CA     \\
            CA^2   \\
            \vdots \\
            CA^{N-1}
        \end{bmatrix}
    \end{equation}
    是列满秩的，即$\rank\mathcal{O}=N$。
\end{theorem}

以上两个定理的证明都与\refsub{state_variable_analysis}中一致。

最后，我们使用单边z变换来分析由状态变量方程描述的系统。对方程组
\begin{equation}
    \begin{aligned}
        \mathbf{q}[n+1]&=A\mathbf{q}[n]+B\mathbf{e}[n], \\
        \mathbf{r}[n]&=C\mathbf{q}[n]+D\mathbf{e}[n], \\
        \mathbf{q}[0]&=\mathbf{q}_0,
    \end{aligned}
\end{equation}
两边同时乘以$u[n]$并做单边z变换，有
\begin{equation}
    \begin{aligned}
        z\mathbf{Q}(z)-z\mathbf{q}_0&=A\mathbf{Q}(z)+B\mathbf{E}(z), \\
        \mathbf{R}(z)&=C\mathbf{Q}(z)+D\mathbf{E}(z).
    \end{aligned}
\end{equation}
由前一方程，
\begin{equation}
    \mathbf{Q}(z)=z(zI-A)^{-1}\mathbf{q}_0+(zI-A)^{-1}B\mathbf{E}(z)=z\Phi(z)\mathbf{q}_0+\Phi(z)B\mathbf{E}(z),
\end{equation}
代入后一方程，得到
\begin{align}
    \mathbf{R}(z)&=C\sbr{z\Phi(z)\mathbf{q}_0+\Phi(z)B\mathbf{E}(z)}+D\mathbf{E}(z) \nonumber \\
    &=zC\Phi(z)\mathbf{q}_0+[C\Phi(z)B+D]\mathbf{E}(z).
\end{align}
因此系统函数为
\begin{equation}
    H(z)=C\Phi(z)B+D.
\end{equation}

对于一个$q$输入、$p$输出的线性时不变系统，其系统函数为$p\times q$的矩阵值函数。特别地，对于单输入单输出系统，系统函数为一个分式函数。

利用系统函数同样可以判断系统的可控性与可观性：

\begin{claim}
    设单输入单输出系统的状态变量方程为
    \begin{equation}
        \mathbf{q}[n+1]=A\mathbf{q}[n]+B\mathbf{e}[n],\quad \mathbf{r}[n]=C\mathbf{q}[n],
    \end{equation}
    则系统的系统函数为
    \begin{equation}
        H(z)=C\Phi(z)B.
    \end{equation}
    如果在系统函数的计算过程中，没有出现零极点的抵消，即系统函数的分子和分母没有公共因子，则系统是可控且可观的，否则系统不完全可控或不完全可观。
\end{claim}
